% Section 5: Recommandations générales
% Fichier: sections/05_recommandations.tex

\chapter{Recommandations générales}

\section{Contre-mesures correctives}

\subsection{Injections}
\begin{itemize}
    \item Utiliser systématiquement des requêtes paramétrées
    \item Éviter la concaténation de chaînes pour construire des requêtes SQL
    \item Valider et nettoyer toutes les entrées utilisateur
\end{itemize}

\subsection{XSS}
\begin{itemize}
    \item Échapper les données utilisateur selon le contexte (HTML, JavaScript, CSS, URL)
    \item Implémenter une Content Security Policy (CSP) stricte
    \item Utiliser des frameworks qui échappent automatiquement les données (Angular, React)
\end{itemize}

\subsection{Authentification}
\begin{itemize}
    \item Implémenter un verrouillage de compte après plusieurs tentatives échouées
    \item Utiliser l'authentification à deux facteurs (2FA)
    \item Stocker les mots de passe avec un algorithme de hachage fort (bcrypt, Argon2)
\end{itemize}

\subsection{Contrôle d'accès}
\begin{itemize}
    \item Vérifier les autorisations pour chaque requête
    \item Utiliser des UUID au lieu d'identifiants séquentiels
    \item Implémenter le principe du moindre privilège
\end{itemize}

\subsection{Configuration}
\begin{itemize}
    \item Désactiver les messages d'erreur détaillés en production
    \item Supprimer les fonctionnalités de débogage
    \item Mettre à jour régulièrement les dépendances
\end{itemize}

\section{Recommandations prioritaires}
\begin{priorityrecommendation}{Critique}
    Corriger immédiatement toutes les vulnérabilités permettant un accès non autorisé ou une exécution de code à distance.
\end{priorityrecommendation}

\begin{priorityrecommendation}{Élevée}
    Mettre en place une politique de sécurité des contenus (CSP) et valider toutes les entrées utilisateur.
\end{priorityrecommendation}

\begin{priorityrecommendation}{Moyenne}
    Implémenter un mécanisme de verrouillage de compte et utiliser des hachages forts pour les mots de passe.
\end{priorityrecommendation}
